
When a medical imaging dataset shifts from X-ray to MRI scans without formal notice, models trained on deprecated data silently degrade. This failure mode is systemic: major machine learning dataset repositories lack the automated deprecation mechanisms standard in software package managers for decades. HuggingFace Datasets Hub hosts over 60,000 datasets with informal versioning practices where maintainers manually track deprecation through README updates and GitHub discussions. This creates information asymmetry: users unknowingly train on stale data while maintainers lack tools to measure whether deprecation notices reach affected users.

Dataset deprecation affects reproducibility, model maintenance, and deployment reliability---core infrastructure problems that scale with machine learning adoption. Current practice relies on documentation-only signals (README warnings, dataset card annotations) that fail at three critical infrastructure points: (1) automated detection of deprecation candidates via usage and maintenance signals, (2) task-specific successor mapping that captures context-dependent replacement paths, and (3) point-of-use adoption tracking that quantifies whether users migrate to recommended successors. Software package managers (NPM, PyPI, Maven) solved analogous problems for code dependencies via automated deprecation warnings and version constraints, but machine learning datasets have been treated as static artifacts rather than living dependencies requiring lifecycle management.

Datasets are dependencies with deprecation lifecycles similar to software packages but with unique requirements. Health metrics must use usage velocity rather than semantic versioning. Successor graphs must be context-aware because task matters: ImageNet may transition to ImageNet-v2 for robustness evaluation but to ImageNet-21k for pretraining. Load-time instrumentation must track adoption events to quantify deprecation mechanism efficacy. No existing system combines automated health metrics, context-aware successor graphs, and instrumented policies in an integrated infrastructure.

This work presents a three-component formal deprecation system validated through mechanistic experiments on proof-of-concept scale infrastructure. The contributions are:

\begin{enumerate}
\item \textbf{System Design:} Three-component architecture integrating automated health metrics (velocity, emergence, issue signals), context-aware successor graphs (task-conditional replacement paths), and load-time instrumentation (adoption measurement telemetry).
\end{enumerate}

\begin{enumerate}
\item \textbf{Mechanistic Validation:} Five sub-hypothesis experiments (h-e1, h-m1, h-m2, h-m3, h-m4) demonstrating feasibility on simulated/mock data. Health metrics achieved 88.3\% precision (target $\geq 60$\%), context inference achieved 100\% accuracy on synthetic patterns (target $\geq 70$\%), and instrumentation added <5\% overhead (target <10\%).
\end{enumerate}

\begin{enumerate}
\item \textbf{Measurement Infrastructure:} First quantitative telemetry infrastructure for dataset adoption tracking with 100\% capture rate (95\% CI: 99.05\%-100\%), enabling future efficacy studies comparing formal versus informal deprecation mechanisms.
\end{enumerate}

Validated at proof-of-concept scale on 100-1000 samples with simulated HuggingFace-compatible data, results demonstrate that formal dataset deprecation mechanisms are mechanistically feasible. Efficacy measurement (adoption lift versus baseline) requires deployment-scale testing and was not performed in this study.
